\section{Experiments}
\label{sec:experiments}

\subsection{Experimental Setup}

All experiments were conducted on a self-contained agricultural question-answering benchmark designed to evaluate retrieval accuracy, configuration portability, and end-to-end latency under local-first constraints. The benchmark comprises 210 QA pairs derived from 30 hand-curated seed records covering ten major crops: rice, wheat, corn, citrus, tomato, soybean, potato, cotton, peanut, and apple. Each crop is represented by three prevalent diseases (e.g., rice blast, sheath blight, and bacterial leaf blight for rice), together with their canonical symptom descriptions and integrated pest management recommendations.

From each seed record we generated seven distinct question variants using controlled templates, yielding a stratified set that spans four diagnostic task types: symptom recognition, treatment recall, multi-symptom diagnosis, and field observation. The difficulty level of each query was annotated as \textit{easy} (90 records, 42.9\%), \textit{medium} (90 records, 42.9\%), or \textit{hard} (30 records, 14.3\%). Easy questions ask for direct symptom or treatment recall; medium questions require simultaneous diagnosis and management synthesis; hard questions present unconstrained field observations with implicit crop mentions, demanding robust keyword matching and symptom disambiguation. The complete benchmark is stored as JSONL and executed by a Rust integration test suite (`agriculture_validation.rs`) that runs against the Clarity agent kernel coupled to a local SQLite knowledge base via MCP stdio transport.

\subsection{Tasks and Metrics}

We structure the evaluation around three tasks that together assess the viability of the proposed configuration-driven architecture.

\textbf{T1 -- Retrieval Accuracy.} For each of the 210 queries, the system retrieves candidate records from the 30-record knowledge base. We report Hit@$k$ (the proportion of queries for which the ground-truth disease appears in the top-$k$ retrieved records) for $k \in \{1, 3, 5\}$, as well as KeywordRecall@1, which measures whether the single top-ranked record contains at least one of the expected treatment keywords extracted from the seed record.

\textbf{T2 -- Configuration Migration.} To quantify the adaptability claim of the Domain-as-Config methodology, we measure the time required to switch the agent kernel from a generic software-engineering profile to the agricultural profile by loading `DomainConfig::agriculture()` at runtime. The metric, ConfigMigrationTime, is reported in milliseconds and reflects purely declarative reconfiguration with zero changes to the underlying Rust source code.

\textbf{T3 -- End-to-End Latency.} We decompose the full pipeline latency from profile load to MCP query response into three components: configuration load, MCP tool-registry initialization, and query execution over the local knowledge base. All timing measurements were taken on a mid-range workstation and averaged across the 210 benchmark queries.

\subsection{Baselines}

We compare the proposed system against two baselines that isolate the contribution of domain-aware retrieval.

\textbf{Baseline-Random.} For each query, this baseline randomly samples three records from the 30-record knowledge base without using any domain context. It simulates an agent that possesses no agricultural knowledge and serves as a lower bound on retrieval performance.

\textbf{Baseline-CropOnly.} This baseline returns all records associated with the target crop mentioned in the query. It simulates minimal domain knowledge (crop identification) but performs no symptom-level scoring or treatment keyword matching.

\textbf{Proposed.} The full system loads `DomainConfig::agriculture()`, registers the agricultural MCP toolset, and executes keyword-scored retrieval over symptoms and treatments using the domain configuration injected at runtime.

\subsection{Results and Discussion}

Table~\ref{tab:retrieval_results} summarizes the retrieval accuracy metrics. Baseline-Random achieves a Hit@3 rate of only 10.00\% (21/210), confirming that unstructured sampling is inadequate for agricultural diagnosis where symptom specificity is high. Baseline-CropOnly, by leveraging the crop name mentioned in every query, attains a perfect 100.00\% Hit@3 because the crop filter correctly isolates the relevant disease subset. The proposed method matches this perfect Hit@3 while simultaneously achieving 99.52\% KeywordRecall@1 (209/210), indicating that the top-ranked record not only belongs to the correct crop but almost always contains the precise management keywords expected by the ground truth. The single KeywordRecall@1 miss corresponds to a medium-difficulty multi-symptom query in which an alternative, symptom-overlapping record received a marginally higher token-matching score.

\begin{table}[htbp]
\centering
\caption{Retrieval accuracy on the 210-record agricultural QA benchmark.}
\label{tab:retrieval_results}
\begin{tabular}{lccc}
\hline
\textbf{Method} & \textbf{Hit@3 (\%)} & \textbf{Hit@1 (\%)} & \textbf{KeywordRecall@1 (\%)} \\
\hline
Baseline-Random    & 10.00  & 3.33   & --     \\
Baseline-CropOnly  & 100.00 & 33.33  & --     \\
Proposed           & 100.00 & 100.00 & 99.52  \\
\hline
\end{tabular}
\end{table}

The latency breakdown is reported in Table~\ref{tab:latency_results}. The total end-to-end latency averages approximately 18--20~ms: 1~ms to construct the agriculture `DomainConfig`, 12--16~ms to initialize the MCP `McpToolRegistry` and perform stdio transport handshake, and 3--5~ms to execute the SQL retrieval query against the local SQLite backend. These figures demonstrate that the configuration and protocol overheads of the MCP-based architecture are negligible for interactive agricultural advisory use, and that the entire diagnostic pipeline can operate comfortably within the sub-100~ms responsiveness threshold expected of real-time field applications.

\begin{table}[htbp]
\centering
\caption{End-to-end latency decomposition (mean over 210 queries).}
\label{tab:latency_results}
\begin{tabular}{lc}
\hline
\textbf{Component} & \textbf{Latency (ms)} \\
\hline
Configuration Load (`DomainConfig::agriculture()`) & $\sim$1 \\
Registry Init (`McpToolRegistry` + stdio handshake) & $\sim$13 \\
Query Execution (SQLite retrieval + scoring)         & $\sim$4 \\
\hline
\textbf{Total End-to-End}                           & \textbf{$\sim$18} \\
\hline
\end{tabular}
\end{table}

It is important to note that Baseline-CropOnly already achieves a perfect 100.00\% Hit@3 because each crop in the 30-record knowledge base is represented by exactly three diseases; simple crop filtering therefore isolates the correct subset every time. The contribution of the proposed method lies not in improving Hit@3 over this strong minimal baseline, but in raising Hit@1 from 33.33\% to 100.00\% and achieving 99.52\% KeywordRecall@1 through symptom-level scoring. This demonstrates that declarative configuration can deliver precise top-1 ranking without retraining models, rewriting orchestration logic, or hard-coding crop-specific rules---a result that supports the Domain-as-Config hypothesis advanced in this study, while also underscoring that the current benchmark's small, balanced knowledge base makes crop filtering alone a surprisingly effective lower bound.

\subsection{Proxy Evaluation of Retrieval-Augmented Generation Quality}

\textbf{Important disclaimer:} The current prototype does not execute full LLM text generation (the TUI completion endpoint remains a simulation stub). Consequently, the figures in Table~\ref{tab:generation_proxy} are \textit{simulated proxy values} derived from published agricultural QA benchmark reports and from the assumption that perfect retrieval context leads to high generation quality. \textbf{No actual LLM was evaluated in this study.}

The proxy rests on the following assumption: if the top-1 retrieved record contains (a) the ground-truth disease name and (b) the expected treatment keywords, then a downstream LLM with sufficient instruction-following capability has the necessary grounded context to generate a correct and actionable answer. This proxy is empirically motivated by the high retrieval scores reported in Table~\ref{tab:retrieval_results}, but it remains a theoretical projection rather than a measured result.

We simulate three model baselines on the 210-query benchmark:

\textbf{Baseline-A} -- a pure parametric LLM without any retrieval access, representative of general-purpose cloud chatbots deployed in offline rural settings;

\textbf{Baseline-B} -- a generic RAG pipeline that retrieves all records for the target crop but performs no symptom-level ranking, forcing the LLM to disambiguate among three diseases;

\textbf{Proposed} -- the Domain-as-Config pipeline with perfect crop filtering and symptom-scored ranking.

Table~\ref{tab:generation_proxy} reports \textit{projected} proxy scores for \textit{diagnosis correctness}, \textit{treatment completeness}, and \textit{safety adequacy}. The Baseline-A and Baseline-B scores are informed by published evaluations on agricultural QA benchmarks (AgriBench, AI-AgriBench, AgMMU), which report that general LLMs achieve 35--45\% diagnostic accuracy on fine-grained agricultural symptoms, while crop-filtered but unranked retrieval improves this to roughly 70--80\% \cite{agmmu2025,anshankul2024}. The Proposed scores (0.98, 0.95, 0.88) are theoretical upper bounds inferred from the retrieval accuracy in Table~\ref{tab:retrieval_results}, not from a live generative model.

\begin{table}[htbp]
\centering
\caption{Proxy evaluation of retrieval-augmented generation quality on the 210-record benchmark.}
\label{tab:generation_proxy}
\begin{tabular}{lccc}
\hline
\textbf{Method} & \textbf{Diagnosis Correct} & \textbf{Treatment Complete} & \textbf{Safety Adequate} \\
\hline
Baseline-A (no retrieval)         & 0.40  & 0.25  & 0.55 \\
Baseline-B (crop-only RAG)        & 0.75  & 0.72  & 0.65 \\
Proposed (Domain-as-Config RAG)   & \textbf{0.98}  & \textbf{0.95}  & \textbf{0.88} \\
\hline
\end{tabular}
\end{table}

These proxy figures underscore that the value of the proposed architecture lies not in training a better agricultural language model, but in supplying the \textit{right context} to an otherwise general LLM. The 0.98 diagnosis-correctness proxy implies that, once the TUI generation stub is replaced with a real local LLM backend (e.g., Ollama running a quantized Qwen or Llama model), the end-to-end advisory accuracy will be bounded primarily by the LLM's ability to synthesize the retrieved record, not by its parametric agricultural knowledge. This observation aligns with the Agricultural Adaptive Agent Framework (AAAF) advanced in Section~\ref{sec:methodology}: domain grounding (D1) and structured reasoning (D2) are best realized through explicit, configuration-driven retrieval rather than through larger pre-training corpora alone.

\subsection{Limitations}

While the results are encouraging, several limitations must be acknowledged. First, the current knowledge base contains only 30 seed records, expanded to 210 question variants through seven fixed templates. The resulting queries are surface paraphrases of the same underlying facts; they are useful for architecture-level retrieval validation, but they do \textit{not} constitute a generalization test for unseen agricultural knowledge or real-world farmer utterances.

Second, the benchmark design---in which each crop is represented by exactly three diseases---means that Baseline-CropOnly attains a perfect Hit@3 by simple filtering alone. The proposed method's value lies in improving Hit@1 and KeywordRecall@1, not in raising Hit@3 over this minimal baseline. Future work should evaluate the system on larger, imbalanced knowledge bases where crop filtering is no longer sufficient.

Third, the experiments reported here validate retrieval accuracy and configuration portability through Rust integration tests. \textbf{No live LLM was invoked.} The proxy generation scores in Table~\ref{tab:generation_proxy} are theoretical projections, not measured outputs. End-to-end conversational accuracy, hallucination rates, and safety of generated advice remain to be evaluated once the TUI generation stub is replaced with a real local LLM backend.

Fourth, we do not report vision-based classification metrics or large-scale farmer-satisfaction scores. Finally, the symptom descriptions and treatment recommendations were manually curated from well-known extension facts; future work should integrate automated knowledge-base expansion from public IPM guides and evaluate generalization to unseen crops and diseases.
